\scanpage{291}
% archive sheet number: 292
b) $\varphi_N:\mathcal M_N\longrightarrow\mathcal M$
\[
\left|\begin{aligned}
(\varphi_N(Q))_i&=\operatorname{Hom}_{\mathcal M}^{i}(\tau_{N-i}(S),\overset{\substack{\inserted{i_N Q}\\\Vert}}{Q})(\longrightarrow\operatorname{Hom}_{\mathcal M}^{0}(S,Q)=Q_{\unsure{i}}\unsure{)}\\
&=\operatorname{Hom}_{\mathcal M}^{\deleted{\unsure{i|*}}}(\tau_{N-i}(S)[-i],Q)
\end{aligned}\right.\tag*{(10.2)}
\]
Avec \deleted{\unsure{leur}} structure de module provenant (pour $k\in\mathbb Z,\ u\in S_{\inserted{k}}^{\deleted{\unsure{k|*}}}$) des opérations, multiplication par $u$ :
\[
\tau_{N-i}(S)\xleftarrow{\deg k}\tau_{N-i-k}(S)
\]
en transposant ; on a un hom \deleted{dans} \deleted{\unsure{mod}} $\mathcal M$ décrit dans (10.2)
\[
\varphi_N(Q)\longrightarrow i_N(Q)\tag*{(10.2.1)}
\]
d'où en appliquant $\tau_N$
\[
\tau_N\varphi_N Q\xrightarrow[\text{adj.}]{\sim}Q\tag*{(10.2.2)}
\]
qui est un \unsure{isom|isom.} (c'est un hom d'adjonction)

L'autre hom d'adjonction
\[
\deleted{\varphi_N\tau_N P\longleftarrow}P\longrightarrow\varphi_N\tau_N P
\]
est déduit des \deleted{\unsure{h'}\ \unsure{d'adjonction}} \inserted{hom}
\[
\deleted{\tau_N P\longrightarrow P}\quad \operatorname{Hom}_{\unsure{\mathcal M}}^{i}(\tau_{N\deleted{\unsure{-i}}}S,P)
\]
\noindent (10.2.3)
\diagram{\begin{tikzcd}[ampersand replacement=\&,column sep=small,row sep=large]
P_i\arrow[rr]\arrow[dr,"\unsure{\sim|2|\simeq}"']\&\&(\varphi_N\tau_N P)_i\simeq\operatorname{Hom}^{i}(\tau_{N-i}S,\tau_N P)\\
\&\operatorname{Hom}^{i}(S,P)\arrow[ur]\&
\end{tikzcd}}
\margin{tout hom de degré $i$ \deleted{de $S$ dans $P$} $\tau_{N-i}S$ dans $P$ est à valeurs dans $\tau_{N\deleted{\unsure{-i}}}P\subset P$}

c) (NB On n'utilisera pas $q_N$, on le met par acquit de conscience) :
\[
q_N(P)=P/\text{sous-module engendré par les }P_i\text{ pour }i<N
\]
